
\maketitle

\begin{abstract}
Building on the canonical probability space representation established
in \emph{Observational Foundations of Probability}, this work studies
the predictive structure that emerges from compatible observable
experiments.

Given a future observable $F:\Omega\to O_F$, prediction induces an
equivalence relation on observable states determined by equality of
conditional laws.  The resulting quotient defines a minimal predictive
query $Q_*$ whose states correspond exactly to distinct predictive
laws of $F$.  A canonical positive linear operator
\[
(K_{Q_*}g)(q_*) = E[g(F)\mid Q_*=q_*]
\]
then acts on functions of this predictive state space, providing an
observable representation of prediction expressed entirely in terms of
conditional laws.  The minimal predictive query is universal: any
other predictive statistic factors through $Q_*$ via the
Doob--Dynkin lemma, making $Q_*$ the coarsest observable
representation retaining the full predictive law of $F$.
\end{abstract}

\paragraph{Relation to the observational foundations.}
\emph{Observational Foundations of Probability} establishes that
compatible observable experiments admit a canonical probability space
$(\Omega,\sigma(\Qcal),P)$.  The present work studies the predictive
structure that emerges once this representation is available:
prediction induces an equivalence on observable states whose quotient
yields a minimal predictive query $Q_*$.  Conditional expectations of
future observables therefore factor through this predictive state
space.

The predictive structure developed in this note proceeds through the
following sequence of constructions.

The logical progression of the present work is:
\[
\begin{aligned}
&\text{canonical experiment } (\Omega,\sigma(\Qcal),P)\\[6pt]
&\downarrow\\[6pt]
&\text{future observable } F:\Omega\to O_F\\[6pt]
&\downarrow\\[6pt]
&\text{predictive kernels } \Pi_Q(q,\cdot)=P(F\in\cdot\mid Q=q)\\[6pt]
&\downarrow\\[6pt]
&\text{predictive equivalence } q\sim q' \iff \Pi_Q(q,\cdot)=\Pi_Q(q',\cdot)\\[6pt]
&\downarrow\\[6pt]
&\text{minimal predictive query } Q_*:\Omega\to O_{Q_*}\\[6pt]
&\downarrow\\[6pt]
&\text{predictive operator } K_{Q_*}g(q_*)=E[g(F)\mid Q_*=q_*]
\end{aligned}
\]

\paragraph{Main theorem.}
The principal result of this work identifies the minimal observable
state space required for prediction.

\begin{theorem}[Predictive state theorem]
Let $F:\Omega\to O_F$ be a future observable.  Predictive equivalence
induces a minimal predictive query $Q_*$ whose observable states
correspond exactly to distinct predictive laws of $F$.  In particular,
conditional expectations of future observables factor through this
predictive state space.
\end{theorem}

\paragraph{Structure of the paper.}
Section~\ref{sec:standing} recalls the observational framework and
the canonical probability space representation.
Section~\ref{sec:query-experiments} introduces query experiments and
the experiment preorder induced by prediction.
Section~\ref{sec:predictive-experiments} introduces predictive kernels
and the kernel-law map.  Sections~\ref{sec:minimal-predictive}
and~\ref{sec:predictive-sufficiency} develop predictive equivalence,
the minimal predictive query, and the associated factorization theorem,
followed by predictive sufficiency.
Sections~\ref{sec:predictive-operators} and~\ref{sec:operator-closure}
introduce the predictive operator representation and establish operator
closure on the minimal predictive query --- the central theoretical
result of the paper.

\section{Standing framework}\label{sec:standing}

Assume the observational construction of
\emph{Observational Foundations of Probability}.

In particular:

\begin{itemize}
\item A query system $(\Qcal,\preceq)$.
\item The canonical experiment
\[
(\Omega,\sigma(\Qcal),P).
\]
\item Evaluation maps
\[
Q:\Omega\to O_Q.
\]
\item A future observable: a measurable map
\[
F:\Omega\to O_F
\]
with $O_F$ standard Borel, representing a quantity whose distribution
is to be predicted from present queries.
\end{itemize}

Observable functions $g:O_F\to\mathbb R$ are taken to be bounded and
measurable throughout.  Boundedness ensures integrability of $g$ under
every predictive law (since predictive laws are probability measures)
and allows conditional expectations and operator compositions to be
handled by standard Bochner integration without additional
integrability hypotheses.

All other results in this note are derived relative to this observational
structure.  The sequel develops the predictive structure that becomes
available once the observational experiment $(\Omega,P)$ has been
constructed.

\begin{remark}[Program context]
The present paper studies the predictive structure that emerges once
the observational representation theorem has produced the canonical
experiment $(\Omega,\sigma(\Qcal),P)$.
\end{remark}

\begin{remark}[Formalization]
The structures used here---query systems, refinement maps, and the
canonical experiment $(\Omega,P)$---correspond to the objects defined in
the Lean formalization of the QuerySystem framework.  In particular,
evaluation maps
\[
Q:\Omega\to O_Q
\]
are the measurable coordinate projections from the projective-limit
space constructed in the extension theorem.

The assumption that $O_F$ is standard Borel is used not only for the
existence of the predictive kernel $\Pi_Q$ (Rokhlin disintegration)
but also for the measurability of the predictive law map $\phi_Q$ into
$\mathcal P(O_F)$: since $\Pi_Q$ is a Markov kernel, the map
$q\mapsto\Pi_Q(q,\cdot)$ is measurable into $\mathcal P(O_F)$ equipped
with the Giry $\sigma$-algebra (which on standard Borel spaces coincides
with the weak convergence $\sigma$-algebra).  All results involving
$\phi_Q$ and $Q_*$ are derived relative to this measurability.
\end{remark}

For every future observable
\[
F:\Omega\to O_F
\]
regular conditional probabilities define predictive kernels
\[
\Pi_Q(q,\cdot)=P(F\in\cdot\mid Q=q)
\]
for each query $Q:\Omega\to O_Q$.  These kernels are compatible with
refinement of queries: whenever $Q_1\preceq Q_2$,
\[
\Pi_{Q_1}=\Pi_{Q_2}\circ\pi^{Q_2}_{Q_1}.
\]
The remainder of the paper develops the predictive consequences of
these kernels, culminating in the minimal predictive query $Q_*$ and
its associated operator representation.

\bigskip

\section[Query experiments]{Query experiments and the information order}\label{sec:query-experiments}

Each query induces a statistical experiment in the sense of Blackwell and
Le Cam.

\begin{definition}[Query experiment]
For a query $Q$, the pushforward
\[
Q:(\Omega,P)\to (O_Q,Q_\#P)
\]
defines a statistical experiment.
\end{definition}

The refinement preorder on queries corresponds to the standard
informativeness order for experiments.

\begin{theorem}[Experiment preorder]
If
\[
Q_1 \preceq Q_2
\]
then the experiment generated by $Q_1$ is a measurable
post-processing of the experiment generated by $Q_2$.  In particular
there exists a measurable map
\[
\pi^{Q_2}_{Q_1}:O_{Q_2}\to O_{Q_1}
\]
such that
\[
Q_1=\pi^{Q_2}_{Q_1}\circ Q_2.
\]
\end{theorem}

Thus the query system $(\Qcal,\preceq)$ may be interpreted as a preorder
of statistical experiments ordered by informativeness.

\begin{remark}
The present work studies the predictive structure induced on the
canonical experiment $(\Omega,P)$ by future observables.
\end{remark}

\begin{remark}
The experiment preorder provides the structural order in which later
notions of predictive sufficiency and minimality will be expressed.
\end{remark}

\bigskip

\section{Predictive experiments}\label{sec:predictive-experiments}

\begin{definition}[Predictive kernel]\label{def:predictive-kernel}
For a query $Q$, the \emph{predictive kernel} is the regular
conditional distribution of $F$ given $Q$:
\[
\Pi_Q(q,A)
=
P(F\in A\mid Q=q),
\qquad
q\in O_Q,\; A\in\mathcal B(O_F).
\]
Since $O_F$ is standard Borel, a regular conditional distribution
exists, so $\Pi_Q$ is a well-defined Markov kernel
\[
\Pi_Q:O_Q\times\mathcal B(O_F)\to[0,1].
\]
The kernel determines the \emph{predictive law map}
\[
\phi_Q:O_Q\to \Pcal(O_F),
\qquad
\phi_Q(q)=\Pi_Q(q,\cdot).
\]
Measurability of $\phi_Q$ follows from the fact that $\Pi_Q$ is a
kernel: for each $A\in\mathcal B(O_F)$ the map $q\mapsto \Pi_Q(q,A)$
is measurable, so $\phi_Q$ is measurable into $\Pcal(O_F)$ equipped
with the weak convergence $\sigma$-algebra.
\end{definition}

\begin{remark}[Regular conditional distribution]
The kernel $\Pi_Q$ satisfies
\[
\Pi_Q(Q(\omega),A)
=
P(F\in A\mid Q)(\omega)
\qquad \text{$P$-a.s.}
\]
It is defined $Q_\#P$-almost surely in $q$.
\end{remark}

\begin{proposition}[Predictive compatibility]\label{prop:pred-compat}
If $Q_1\preceq Q_2$ then
\[
\Pi_{Q_1}
=
\Pi_{Q_2}\circ\pi^{Q_2}_{Q_1}
\qquad Q_1{}_\#P\text{-a.s.}
\]
\end{proposition}

\begin{proof}
By the tower property of conditional expectation,
for any $A\in\mathcal B(O_F)$,
\[
P(F\in A\mid Q_1)
=
E\bigl[P(F\in A\mid Q_2)\bigm| Q_1\bigr]
=
E\bigl[\Pi_{Q_2}(Q_2,A)\bigm| Q_1\bigr].
\]
Since $Q_1=\pi^{Q_2}_{Q_1}\circ Q_2$ and $\Pi_{Q_2}(Q_2,A)$ is
$\sigma(Q_2)$-measurable, the right-hand side equals
$\Pi_{Q_2}(\pi^{Q_2}_{Q_1}(Q_1),A)$,
which equals $(\Pi_{Q_2}\circ\pi^{Q_2}_{Q_1})(Q_1,A)$.
\end{proof}

This expresses that predictive experiments inherit the refinement order on queries.

\begin{remark}[Necessity of the a.e.\ qualifier]
The equality $\Pi_{Q_1} = \Pi_{Q_2}\circ\pi^{Q_2}_{Q_1}$ holds
$Q_{1\#}P$-almost surely, not pointwise.  This is unavoidable:
the regular conditional distribution $\Pi_Q(q,\cdot)$ is only
determined $Q_\#P$-almost surely in $q$, since two versions of a
conditional distribution may differ on a set of $Q_\#P$-measure
zero.  Any two regular conditional distributions of $F$ given $Q$
therefore agree $Q_\#P$-a.s., and the compatibility identity can
only be asserted in this almost-sure sense.
\end{remark}

\bigskip

\begin{remark}[Predictive statistics]
From the perspective of statistical experiment theory, predictive
queries behave like statistics of the present relative to the future.
The minimal predictive query introduced later will therefore play the
role of a dynamical sufficient statistic: the smallest observable
representation of the present retaining exactly the predictive
information relevant to the future observable $F$.
\end{remark}

The next step is to identify when two observable states determine the
same predictive law.  This leads to the notion of predictive
equivalence and to the construction of the minimal predictive query.

\bigskip

\section{Minimal predictive statistics}\label{sec:minimal-predictive}

\begin{definition}[Predictive equivalence]\label{def:predictive-equivalence}
Let $Q:\Omega\to O_Q$ be a query.  Two observable states
$q,q'\in O_Q$ are said to be \emph{predictively equivalent} if
\[
\phi_Q(q)=\phi_Q(q'),
\]
or equivalently
\[
\Pi_Q(q,\cdot)=\Pi_Q(q',\cdot).
\]
\end{definition}

Two observable states are predictively equivalent whenever they have
the same image under $\phi_Q$, i.e.\ whenever they induce the same
predictive law of $F$.

\begin{theorem}[Minimal predictive query]\label{thm:minimal-predictive-query}
Assume $O_Q$ and $O_F$ are standard Borel.  Define
\[
Q_* := \phi_Q \circ Q : \Omega \to \Pcal(O_F).
\]
Then $Q_*$ is measurable, and
\[
Q_*(\omega) = P(F\in\cdot\mid Q=Q(\omega)).
\]
Two realizations satisfy $Q_*(\omega)=Q_*(\omega')$ if and only if
they are predictively equivalent.  The state space of $Q_*$ is
\[
O_{Q_*} = \phi_Q(O_Q) \subseteq \Pcal(O_F),
\]
the image of the predictive law map, so each state of $Q_*$
corresponds exactly to a distinct predictive law of $F$.
\end{theorem}

\begin{proof}
Measurability of $Q_*=\phi_Q\circ Q$ follows from measurability of
$\phi_Q$ (Definition~\ref{def:predictive-kernel}) and of $Q$.  The
identification $Q_*(\omega)=P(F\in\cdot\mid Q=Q(\omega))$ is immediate
from the definition of $\phi_Q$.  Two realizations satisfy
$Q_*(\omega)=Q_*(\omega')$ iff $\phi_Q(Q(\omega))=\phi_Q(Q(\omega'))$
iff they are predictively equivalent
(Definition~\ref{def:predictive-equivalence}).  The state space
$O_{Q_*}=\phi_Q(O_Q)$ by definition.
\end{proof}

\begin{remark}[Predictive state space]
The state space $O_{Q_*}=\phi_Q(O_Q)\subseteq\Pcal(O_F)$ is a space
of probability measures: each state is the predictive law of $F$
associated with a present observation.  The minimal predictive query
organizes prediction by mapping each observable state to its
associated predictive law.
\end{remark}

The predictive law of $F$ therefore depends on $\omega$ only through
$Q_*(\omega)$.  The following theorem makes this factorization precise
and establishes its universal minimality.

\begin{theorem}[Predictive factorization theorem]\label{thm:predictive-factorization}
Let $Q_*=\phi_Q\circ Q$ be the minimal predictive query of
Theorem~\ref{thm:minimal-predictive-query}, and let
$\Phi:O_{Q_*}\to\Pcal(O_F)$ be the identity inclusion
$\Phi(q_*)=q_*$.

Then
\[
P(F\in\cdot\mid Q_*=q_*)=\Phi(q_*)=q_*,
\]
so the predictive law $\Pi(\omega)=P(F\in\cdot\mid\omega)$ factors as
\[
\Pi = \Phi\circ Q_*.
\]

Moreover $Q_*$ is universal: if $R:\Omega\to O_R$ is any query and
$\psi:O_R\to\Pcal(O_F)$ is measurable with $\Pi=\psi\circ R$, then
$Q_*$ is measurable with respect to $\sigma(R)$, and there exists a
measurable map $\theta:O_R\to O_{Q_*}$ such that $Q_*=\theta\circ R$.

In particular the following diagram commutes:
\[
\begin{array}{ccc}
\Omega & \xrightarrow{\;\;R\;\;} & O_R \\[4pt]
\downarrow\scriptstyle{Q_*} & & \downarrow\scriptstyle{\psi} \\[4pt]
O_{Q_*} & \xrightarrow{\;\;\Phi\;\;} & \Pcal(O_F)
\end{array}
\]

Thus $Q_*$ is the minimal observable statistic through which the
predictive law of the future experiment factors.
\end{theorem}

\begin{proof}
\textit{Factorization.}
Since $Q_*(\omega)=\phi_Q(Q(\omega))=\Pi_Q(Q(\omega),\cdot)
=P(F\in\cdot\mid Q)(\omega)$, we have
$P(F\in\cdot\mid Q_*=q_*)=q_*=\Phi(q_*)$,
so $\Pi=\Phi\circ Q_*$.

\textit{Universality.}
Suppose $\Pi=\psi\circ R$.  Then $Q_*=\Pi=\psi\circ R$ is
$\sigma(R)$-measurable.  By the Doob--Dynkin lemma applied to the
measurable map $Q_*:\Omega\to\Pcal(O_F)$ and the $\sigma$-algebra
$\sigma(R)$, there exists a measurable $\theta:O_R\to O_{Q_*}$ with
$Q_*=\theta\circ R$.

\textit{Minimality.}
The universality step shows $Q_*$ factors through any other predictive
statistic.  Therefore $Q_*$ is the coarsest measurable function of
$\omega$ carrying the full predictive law.
\end{proof}

\begin{corollary}[Canonical predictive operator]\label{cor:canonical-operator}
For every bounded measurable $g:O_F\to\mathbb R$, define
\[
(K g)(q_*)
=
\int g(f)\,q_*(df)
=
E[g(F)\mid Q_*=q_*].
\]
Then $K:L^\infty(O_F)\to L^\infty(O_{Q_*})$ is a well-defined positive
linear operator.
\end{corollary}

\begin{proof}
Since $q_*\in\Pcal(O_F)$ is a probability measure, integration against
$q_*$ is a positive linear functional.  The function $q_*\mapsto\int g\,dq_*$
is measurable in $q_*$ by measurability of $\phi_Q$ and boundedness
of $g$.  Positivity and linearity of $K$ are inherited from those of
integration.
\end{proof}

\begin{remark}[Predictive operators]
The predictive factorization therefore induces operators acting on
functions of the predictive state.  These operators represent the
evolution of predictive expectations of future observables.

In particular, if $F_t$ denotes the observable at prediction horizon
$t$, the operators
\[
(K_t g)(q_*)=E[g(F_t)\mid Q_*=q_*]
\]
form the predictive evolution operators studied in the next paper.
\end{remark}

\begin{remark}[Predictive horizon quotient]
The canonical experiment space $\Omega$ may be interpreted as an
observable horizon of histories.  Predictive equivalence identifies
two horizon configurations whenever they induce the same conditional
law of the future observable.  The minimal predictive query $Q_*$ is
therefore the quotient of the observable horizon by predictive
indistinguishability.  In this sense $Q_*$ is the minimal observable
boundary state needed for autonomous prediction.
\end{remark}

\begin{remark}[Lattice interpretation]
When the experiment preorder admits infima, $Q_*$ equals the meet of
all predictively sufficient queries:
\[
Q_* = \bigwedge_{Q\in\mathcal S} Q.
\]
Thus $Q_*$ is the largest common coarsening of all sufficient
experiments: it extracts exactly the predictive information shared by
all of them, no more.
\end{remark}

\begin{remark}[Relation to Blackwell sufficiency]
In the language of statistical experiment theory, the minimal predictive
query plays the role of the minimal sufficient experiment for predicting
the future observable $F$: it is the coarsest observable statistic that
retains the full predictive law.
\end{remark}

\begin{remark}[Doob--Dynkin factorization of the future experiment]
The predictive factorization theorem (Theorem~\ref{thm:predictive-factorization})
shows that the minimal predictive query $Q_*$ is the minimal
Doob--Dynkin factor through which the future experiment factors.
In this sense $Q_*$ is simultaneously a minimal sufficient statistic
and a minimal Doob--Dynkin factor for prediction.
\end{remark}

\bigskip

\section{Predictive sufficiency}\label{sec:predictive-sufficiency}

The minimal predictive query constructed in the previous section
captures all predictive information about the future observable.
This section interprets that construction through the notion of
predictive sufficiency and relates it to classical ideas from
statistical experiment theory.

\begin{definition}[Predictive sufficiency]
A query $Q$ is predictively sufficient if for every admissible query $Q'$
there exists a measurable map
\[
\psi:O_Q\to O_{Q'}
\]
such that
\[
\Pi_{Q'}=\Pi_Q\circ\psi.
\]
\end{definition}

Predictively sufficient queries therefore retain all predictive
information about the future observable.

\begin{theorem}[Predictive sufficiency characterization]
A query $Q$ is predictively sufficient if and only if
\[
F \perp Q' \mid Q
\]
for all admissible queries $Q'$.
\end{theorem}

\begin{proof}
($\Rightarrow$)
Suppose $Q$ is predictively sufficient, so for every $Q'$ there exists
measurable $\psi:O_Q\to O_{Q'}$ with $\Pi_{Q'}=\Pi_Q\circ\psi$.
Then for any $A\in\mathcal B(O_F)$,
\[
P(F\in A\mid Q,Q')
=
P(F\in A\mid Q),
\]
which is exactly $F\perp Q'\mid Q$.

($\Leftarrow$)
If $F\perp Q'\mid Q$ for all $Q'$, then for any $A$,
$P(F\in A\mid Q,Q')=P(F\in A\mid Q)$.
In particular $P(F\in A\mid Q')=E[P(F\in A\mid Q)\mid Q']
=E[\Pi_Q(Q,A)\mid Q']$,
so $\Pi_{Q'}=\Pi_Q\circ\psi$ for the conditional expectation map
$\psi(q')=E[Q\mid Q'=q']$ (when $Q$ is a function of $Q'$ this is the
refinement map; in general it is a regular conditional kernel).
Thus $Q$ is predictively sufficient.
\end{proof}

\begin{remark}[Operator formulation and boundedness]
In the operator formulation of Section~\ref{sec:predictive-operators},
predictive sufficiency of $Q$ implies that for every bounded measurable
$g:O_F\to\mathbb R$,
\[
E[g(F)\mid Q] = E[g(F)\mid Q_*]
\qquad P\text{-a.s.}
\]
The Lean formalization of this statement (\texttt{predictive\_sufficiency})
requires an explicit uniform bound on $g$ for the Bochner integration step.
\end{remark}

\section{Predictive operators}\label{sec:predictive-operators}

\begin{definition}[Predictive operator]
For a query $Q$ and bounded measurable function
$g:O_F\to\mathbb R$, define
\[
(K_Q g)(q)
=
E[g(F)\mid Q=q].
\]
\end{definition}

\begin{proposition}[Predictive operator well-definedness]
For bounded measurable $g$, $K_Q$ defines a positive linear operator.
\end{proposition}

\begin{theorem}[Predictive operator representation]
Let $F:\Omega\to O_F$ be a future observable and let $Q_*$ be the
minimal predictive query.  Let $g:O_F\to\mathbb R$ be a bounded
measurable function.  Then the map
\[
(K_{Q_*}g)(q_*)
=
E[g(F)\mid Q_*=q_*]
\]
defines a Markov operator on the predictive state space.

Moreover, the predictive evolution of observable expectations factors
through $Q_*$ and is therefore represented entirely by the operator
$K_{Q_*}$.
\end{theorem}

\begin{remark}[Boundedness in the Lean formalization]
The Lean formalization of this theorem (\texttt{predictive\_factorization}
and \texttt{predictive\_sufficiency}) carries an explicit uniform bound
$\exists\,C,\;\forall f,\;|g(f)|\le C$ as a hypothesis.  This is needed
because the proof uses Bochner integration against a family of probability
measures indexed by $q_*$, and boundedness ensures integrability under
every predictive law without further verification.  Boundedness of $g$ is
listed as a standing assumption in Section~\ref{sec:standing}, so the
hypothesis is always satisfied in practice.
\end{remark}

\begin{remark}[Kernel representation]
The predictive operator is the integral operator induced by the
predictive kernel:

\[
(K_Q g)(q)
=
\int g(f)\,\Pi_Q(q,df).
\]

Thus predictive operators form a class of Markov operators on the
observable state space.  In particular, when $Q=Q_*$, the predictive
operator $K_{Q_*}$ gives the canonical operator representation of the
future experiment on the minimal predictive state space.
\end{remark}

\begin{proposition}[Koopman deterministic specialization]
If the predictive kernel is deterministic,
\[
\Pi_Q(q)=\delta_{\phi(q)},
\]
then
\[
(K_Q g)(q)=g(\phi(q)).
\]
\end{proposition}

Thus classical Koopman operators arise as a special case.

\bigskip

\section{Operator closure}\label{sec:operator-closure}

\begin{corollary}[Observable operator closure]
If $Q_*$ is predictively sufficient then
\[
K_{Q_*}g(q)=E[g(F)\mid Q_*=q]
\]
completely determines the predictive evolution of observable
expectations.
In particular, predictive dynamics may be represented entirely as
an operator acting on the predictive state space $Q_*$.
\end{corollary}

\begin{remark}[Representation of predictive dynamics]
The minimal predictive query $Q_*$ provides an observable state space
for prediction: its states correspond exactly to distinct conditional
laws of the future observable $F$.  The predictive operator $K_{Q_*}$
therefore describes the evolution of predictive expectations on this
state space, yielding an observable representation of predictive
dynamics without reference to latent state variables.
\end{remark}

%%% Delay query material moved to Paper 4 (observational_probability.tex)

\section{Conceptual chain}

The preceding sections establish the following structural progression.

\[
\begin{aligned}
&\text{query system}\\
&\downarrow\\
&\text{probabilistic representation}\\
&\downarrow\\
&\text{query experiments and experiment preorder}\\
&\downarrow\\
&\text{predictive experiments and predictive kernels}\\
&\downarrow\\
&\text{predictive equivalence}\\
&\downarrow\\
&Q_*\ (\text{minimal predictive state space})\\
&\downarrow\\
&K_{Q_*}\ (\text{predictive operator on }O_{Q_*})
\end{aligned}
\]

\begin{remark}[Continuation]
The horizon-indexed family $\{K_t\}_{t\ge0}$, its semigroup structure,
the infinitesimal generator, and Koopman--Perron duality are developed
in \emph{Predictive Operator Theory for Observable Dynamical Systems}.
The constructive route to the predictive state via delay embeddings ---
including the probabilistic sufficiency criterion and the Takens
specialization --- is developed in \emph{Observational Probability from
Delay Queries}.
\end{remark}

\section{Conclusion and related work}

\paragraph{What has been established.}
Starting from the canonical probability space $(\Omega,\sigma(\Qcal),P)$
of \emph{Observational Foundations of Probability}, this paper has
developed the predictive structure that emerges from a future observable
$F:\Omega\to O_F$.  The main results are:

\begin{enumerate}
  \item \emph{Predictive kernels.}  Each query $Q$ induces a
    conditional kernel $\Pi_Q(q,\cdot) = P(F\in\cdot\mid Q=q)$, giving
    prediction a concrete query-level representation.
  \item \emph{Minimal predictive query.}  Predictive equivalence $q\sim q'$
    iff $\Pi_Q(q,\cdot)=\Pi_Q(q',\cdot)$ factors $Q$ through a minimal
    statistic $Q_*$ whose states are in bijection with distinct predictive
    laws of $F$.
  \item \emph{Doob--Dynkin universality.}  Any measurable predictive
    statistic factors through $Q_*$, making the minimal predictive query
    the coarsest observable representation retaining full predictive
    content.
  \item \emph{Predictive operator.}  Conditional expectation
    $(K_{Q_*}g)(q_*) = E[g(F)\mid Q_*=q_*]$ defines a positive linear
    contraction on $O_{Q_*}$ that is closed under the predictive state
    algebra.
\end{enumerate}

\paragraph{Relation to classical sufficiency.}
The minimal predictive query $Q_*$ is an instance of a sufficient
statistic in the sense of Fisher and Halmos--Savage: it retains all
information relevant to predicting $F$ while discarding information
that is irrelevant.  The framework of Blackwell~\cite{Blackwell} and
Le Cam~\cite{LeCam} on comparison and sufficiency of statistical
experiments provides the conceptual background; the present work
specialises their comparison order to the query setting and identifies
the minimal representative in a canonical way via the Doob--Dynkin
lemma.  The experiment preorder induced by prediction is a structural
refinement of the Blackwell order restricted to the observable interface.

\paragraph{Relation to de Finetti and predictive sufficiency.}
In de Finetti's representation theorem~\cite{deFinetti}, the
exchangeable measure is disintegrated over an invariant $\sigma$-field
that plays the role of the predictive state.  The present framework
provides a direct construction of that $\sigma$-field from the
observable query structure, without requiring an exchangeability or
invariance assumption.  Predictive sufficiency (Section~\ref{sec:predictive-sufficiency})
is the query-theoretic counterpart of the Radon--Nikodym disintegration
underlying de Finetti's theorem.

\paragraph{What follows.}
The predictive operator $K_{Q_*}$ established here is a single-horizon
object.  \emph{Predictive Operator Theory for Observable Dynamical
Systems} (Paper~3) promotes this to a horizon-indexed family
$\{K_t\}_{t\ge0}$ and shows it forms a strongly continuous Markov
operator semigroup, from which an infinitesimal generator and predictive
Kolmogorov equation follow.  The constructive route to $Q_*$ via delay
embeddings of a dynamical system --- and the connection to Takens'
embedding theorem --- is the subject of \emph{Observational Probability
from Delay Queries} (Paper~4).

\ifcsname thesismode\endcsname\else
\bibliographystyle{plain}
\bibliography{references}
\fi

\section*{Acknowledgements}

The author thanks Claude (Anthropic) for assistance with mathematical
writing, Lean~4 formalization, and technical discussion throughout the
development of this work.  The author gratefully acknowledges financial
support from Dr.\ Simon Bonner and Dr.\ Matt Davison at the University
of Western Ontario, and from MITACS.

